\SourcePage{348}{353}
\section{Crossing invariance}

The representation of amplitudes of scattering $M_{fi}$ by Feynman integrals reveals their remarkable symmetry, consisting in the following.

Any of the incoming external lines of a Feynman diagram can be regarded (without changing the direction of its arrow) as a particle in the initial or antiparticle in the final state, and each outgoing line as a final particle or initial antiparticle. Simultaneously with the transition from particle to antiparticle there also changes the meaning of the 4-momentum ascribed to the line: $p_i \to -p_f$, where $p_i$ is the 4-momentum of the initial particle (in one channel), and $p_f$ is the 4-momentum of the final (in another channel) particle. The polarisation also changes: the initial wave amplitude $u$ is replaced for a particle (say, electron) by $u = u_s$, and the outgoing $\bar{u}$ by $u^*_s$; for a positron $u = u^*_s$. But the transition from $u$ to $u^*$ and from $\bar{u}$ to $\bar{u}^*$ signifies a change of sign of the projection of spin (or of its helicity).

For a photon, as a truly neutral particle, the change of the meaning of an external line means simply the transition from emission of the photon to its absorption or vice versa: an external photonic line with 4-momentum $k$ corresponds either to absorption of a photon with momentum $k_{\text{abs}} = k$, or to emission of a photon with momentum $k_{\text{em}} = -k$ and with opposite sign of the helicity.

Such a change of meaning of external lines is equivalent to the transition from one crossed channel of the reaction to another. Depending on the channel, only the meaning of the arguments of the function changes: when going from particle to antiparticle the replacement $p_i \to -p_f$ is made. About this property of the amplitude of scattering one speaks as about \textit{crossing symmetry}, or \textit{crossing invariance}.

\SourcePage{349}{354}
In terms of the invariant amplitudes introduced in~\S\,70 and of the kinematic invariants, one can say that these are one and the same functions for all channels, but for each channel their arguments run through values in their own physical region. In other words, the Feynman integrals define the invariant amplitudes as analytic functions; their values in the various physical regions are the analytic continuation of a function given in one of the regions. Since the Feynman integral expressions contain singularities, the invariant amplitudes also have singularities; their values in different physical regions are defined by the appropriate rules of going around these integrals (taking into account the rule of going around the poles---cf.~\S\,75).

Let us stress that crossing invariance is nothing more than a property of the scattering matrix, following from the general requirements of space-time symmetry. Subsequently the equality of amplitudes of processes obtained from one another by the permutation of initial and final states with the replacement of all particles by antiparticles (at unchanged momenta $\mathbf{p}$ of all particles and changed in sign projections of all their moments) is required. This is the requirement of $CPT$-invariance\,\footnote{Let us note that the formal description of transition from one of the indicated reactions to another by means of changing the sign of all 4-momenta on Feynman diagrams corresponds to the meaning of the $CPT$-operation as 4-inversion.}\,\footnote{If this or that channel is forbidden by conservation of 4-momentum, then the probability of the transition automatically vanishes by the $\delta$-function figuring in~(64.5) as the common factor.}. Crossing invariance, however, allows this transformation to be made not for all particles simultaneously but also for any individual particle separately.
